\documentclass[12pt]{article}
\usepackage{amsmath,amssymb,amsthm}
\usepackage[margin=1in]{geometry}

\newtheorem{theorem}{Theorem}
\newtheorem{lemma}[theorem]{Lemma}

\title{Problem 281: Pizza Toppings}
\author{Project Euler Solution}
\date{}

\begin{document}
\maketitle

\section{Problem Statement}

A pizza is cut into $m \cdot n$ equal slices. Let $f(m,n)$ denote the number of distinct distributions of $m$ distinct toppings ($m \ge 2$), each on exactly $n$ slices ($n \ge 1$), where rotations are identified but reflections are distinct.

Find $\displaystyle\sum_{\substack{m \ge 2,\; n \ge 1 \\ f(m,n) \le 10^{15}}} f(m,n)$.

\section{Mathematical Development}

The problem is equivalent to counting necklaces of $N = mn$ beads with $m$ distinct colors, each appearing exactly $n$ times, under the cyclic group $C_N$.

\begin{theorem}[Burnside Counting Formula]
Let $N = mn$. Then
\[
f(m,n) = \frac{1}{mn} \sum_{t \mid n} \phi\!\left(\frac{n}{t}\right) \cdot \frac{(mt)!}{(t!)^m}.
\]
\end{theorem}

\begin{proof}
By Burnside's lemma,
\[
f(m,n) = \frac{1}{N} \sum_{k=0}^{N-1} |\operatorname{Fix}(r^k)|.
\]

\textbf{Step 1 (Fixed-point characterization).} A coloring is fixed by the rotation $r^k$ if and only if it has period dividing $g = \gcd(k,N)$. Such a coloring comprises $N/g$ identical blocks of length~$g$, each containing $g/m$ copies of each color. This requires $m \mid g$; otherwise $|\operatorname{Fix}(r^k)| = 0$.

\textbf{Step 2 (Block counting).} When $m \mid g$, the number of valid blocks of length $g$ is the multinomial coefficient
\[
\frac{g!}{\bigl((g/m)!\bigr)^m}.
\]

\textbf{Step 3 (Grouping by gcd).} For each divisor $g$ of $N$, exactly $\phi(N/g)$ values of $k \in \{0,\ldots,N-1\}$ satisfy $\gcd(k,N) = g$. Hence
\[
f(m,n) = \frac{1}{N} \sum_{\substack{g \mid N \\ m \mid g}} \phi\!\left(\frac{N}{g}\right) \cdot \frac{g!}{\bigl((g/m)!\bigr)^m}.
\]

\textbf{Step 4 (Substitution $g = mt$).} Setting $g = mt$ with $t \mid n$, and noting $N/g = n/t$, $(g/m)! = t!$:
\[
f(m,n) = \frac{1}{mn} \sum_{t \mid n} \phi\!\left(\frac{n}{t}\right) \cdot \frac{(mt)!}{(t!)^m}. \qedhere
\]
\end{proof}

\begin{lemma}[Verification]
The formula reproduces the given values: $f(2,1)=1$, $f(2,2)=2$, $f(3,1)=2$, $f(3,2)=16$.
\end{lemma}

\begin{proof}
Direct computation:
\begin{align*}
f(2,1) &= \tfrac{1}{2}\,\phi(1)\,\tfrac{2!}{(1!)^2} = 1, \\
f(2,2) &= \tfrac{1}{4}\bigl[\phi(2)\,\tfrac{2!}{(1!)^2} + \phi(1)\,\tfrac{4!}{(2!)^2}\bigr] = \tfrac{1}{4}[2+6] = 2, \\
f(3,1) &= \tfrac{1}{3}\,\phi(1)\,\tfrac{3!}{(1!)^3} = 2, \\
f(3,2) &= \tfrac{1}{6}\bigl[\phi(2)\,\tfrac{3!}{(1!)^3} + \phi(1)\,\tfrac{6!}{(2!)^3}\bigr] = \tfrac{1}{6}[6+90] = 16. \qedhere
\end{align*}
\end{proof}

\begin{lemma}[Search Bounds]
It suffices to search $2 \le m \le 18$ and $1 \le n \le 29$.
\end{lemma}

\begin{proof}
Since $f(m,1) = (m-1)!$ and $18! > 10^{15}$, we need $m \le 18$. For fixed $m$, $f(m,n)$ grows super-exponentially in $n$, so numerical evaluation confirms $n \le 29$ suffices.
\end{proof}

\section{Algorithm}

\begin{verbatim}
function solve():
    total = 0
    for m = 2 to 18:
        for n = 1 to 29:
            val = f(m, n)
            if val <= 10^15: total += val
    return total

function f(m, n):
    result = 0
    for each divisor t of n:
        result += euler_phi(n/t) * (m*t)! / (t!)^m
    return result / (m*n)
\end{verbatim}

\section{Complexity Analysis}

\begin{itemize}
\item \textbf{Time:} $O(1)$, as all parameters are bounded by small constants (at most $493$ pairs, each with $O(1)$ divisor iterations).
\item \textbf{Space:} $O(1)$.
\end{itemize}

\section{Answer}

\[
\boxed{1485776387445623}
\]

\end{document}
